% BHU PYQ -- Manually transcribed & error-corrected
% Paper: BCHI-326: Industrial Relations and Labour Laws
% Exam: B.Com. Semester VI Examination, 2015-16

\documentclass[11pt,a4paper]{article}
\usepackage[a4paper,top=2.5cm,bottom=2.5cm,left=2.8cm,right=2.8cm,headheight=14pt]{geometry}
\usepackage{amsmath,amssymb}
\usepackage[T1]{fontenc}
\usepackage[utf8]{inputenc}
\usepackage{lmodern,microtype,fancyhdr,enumitem,hyperref,lastpage,parskip}

\hypersetup{
    colorlinks=true,
    urlcolor=black,
    linkcolor=black,
    pdftitle={BCHI-326: Industrial Relations and Labour Laws -- BHU B.Com. Sem VI 2015-16}
}

\pagestyle{fancy}
\fancyhf{}
\renewcommand{\headrulewidth}{0.4pt}
\renewcommand{\footrulewidth}{0.4pt}
\fancyhead[L]{\small   \textit{Banaras Hindu University}}
\fancyhead[R]{\small   \textit{BCHI-326: Industrial Relations and Labour Laws}}
\fancyfoot[C]{\small Page  \thepage\ of \pageref{LastPage}}


\newlist{parts}{enumerate}{2}
\setlist[parts,1]{label=(\alph*),leftmargin=*,itemsep=5pt,topsep=3pt}
\setlist[parts,2]{label=(\roman*),leftmargin=*,itemsep=3pt,topsep=2pt}

\newcommand{\pts}[1]{\hfill   \textit{[#1]}}


\begin{document}


\begin{center}
  {\large\bfseries BANARAS HINDU UNIVERSITY}\\[2pt]
  {\normalsize B.Com. --- Semester VI Examination, 2015-16}\\[6pt]
  \rule{0.8\linewidth}{0.5pt}\\[6pt]
  {\large\bfseries Paper No. BCHI-326: Industrial Relations and Labour Laws}\\[4pt]
  {\normalsize Time: 3 Hours\hfill Maximum Marks: 70}
\end{center}

\rule{\linewidth}{0.4pt}

\smallskip



\noindent   \textit{Instructions: Attempt all questions. All questions carry equal marks (14 marks each).}

\bigskip




\noindent   \textbf{Question 1.} \\
Define Industrial Relations. Comment on the emerging trends in industrial relations in India.\pts{14}

\centerline{   \textbf{OR}}

Discuss the preconditions for a sound industrial relations system and describe the importance of industrial relations.\pts{14}

\medskip\rule{\linewidth}{0.2pt}




\noindent   \textbf{Question 2.} \\
``Politicisation of trade union leads to inter-union rivalry in India.'' Examine this statement.\pts{14}

\centerline{   \textbf{OR}}

What is the concept of a trade union under the Trade Unions Act, 1926? Explain the functions of a trade union.\pts{14}

\medskip\rule{\linewidth}{0.2pt}




\noindent   \textbf{Question 3.} \\
``Prevention of disputes is good for industrial peace.'' Examine this statement and explain the machinery for the prevention of industrial disputes.\pts{14}

\centerline{   \textbf{OR}}

Discuss the main features and provisions of the Industrial Disputes Act, 1947.\pts{14}

\medskip\rule{\linewidth}{0.2pt}

\pagebreak




\noindent   \textbf{Question 4.} \\
Explain the scope of collective bargaining. What are the requirements of an effective collective bargaining process?\pts{14}

\centerline{   \textbf{OR}}

What is the process of collective bargaining? Review the practice of collective bargaining in India.\pts{14}

\medskip\rule{\linewidth}{0.2pt}




\noindent   \textbf{Question 5.} \\
Explain the main provisions of the Factories Act, 1948 regarding the safety and welfare of the workers in a factory.\pts{14}

\centerline{   \textbf{OR}}

Explain the main provisions of the Payment of Wages Act, 1936.\pts{14}

\vfill
\rule{\linewidth}{0.4pt}

\begin{center}\small
\href{https://bhu-exam.vercel.app/}{Click here to visit our website}\\[4pt]
\href{https://me-aryan.vercel.app/}{Click here to visit our contributor}\\[4pt]
\href{https://quashan-me.vercel.app/}{Click here to visit our contributor}
\end{center}

\end{document}
